% SPEC-DB-002 --- Typed Repositories
% Shared preamble for all spec documents.
% Include from each spec: % Shared preamble for all spec documents.
% Include from each spec: % Shared preamble for all spec documents.
% Include from each spec: \input{../preamble.tex} (adjust depth).

\documentclass[11pt,a4paper]{article}

\usepackage[utf8]{inputenc}
\usepackage[T1]{fontenc}
\usepackage{lmodern}
\usepackage[a4paper,margin=2.3cm]{geometry}
\usepackage{parskip}
\usepackage{microtype}
\usepackage[dvipsnames]{xcolor}
\usepackage{enumitem}
\usepackage{listings}
\usepackage{tcolorbox}
\usepackage{titlesec}
\usepackage{fancyhdr}
\usepackage{array}
\usepackage{longtable}
\usepackage{booktabs}
\usepackage{amsmath}
\usepackage{amssymb}
\usepackage{hyperref}
\usepackage{cleveref}

\tcbuselibrary{skins,breakable}

\definecolor{specBlue}{RGB}{30,64,132}
\definecolor{specGray}{RGB}{90,90,90}
\definecolor{specAccent}{RGB}{198,40,40}
\definecolor{specGreen}{RGB}{30,110,60}

\hypersetup{
  colorlinks=true,
  linkcolor=specBlue,
  urlcolor=specBlue,
  citecolor=specBlue,
  pdfborder={0 0 0}
}

\titleformat{\section}{\Large\bfseries\color{specBlue}}{\thesection}{0.75em}{}
\titleformat{\subsection}{\large\bfseries\color{specBlue!80}}{\thesubsection}{0.6em}{}

\makeatletter
\newcommand{\specID}[1]{\def\@specID{#1}}
\newcommand{\specTitle}[1]{\def\@specTitle{#1}}
\newcommand{\specStatus}[1]{\def\@specStatus{#1}}
\newcommand{\specVersion}[1]{\def\@specVersion{#1}}
\newcommand{\specOwner}[1]{\def\@specOwner{#1}}
\newcommand{\specTraces}[1]{\def\@specTraces{#1}}

\specID{SPEC-XXX}
\specTitle{Untitled Spec}
\specStatus{DRAFT}
\specVersion{0.1}
\specOwner{Jeremie Nunez}
\specTraces{}

\newcommand{\makespecheader}{%
  \begin{center}
    {\LARGE\bfseries\color{specBlue}\@specTitle}\\[0.4em]
    {\small\color{specGray}\texttt{\@specID} \quad v\@specVersion \quad \textsc{\@specStatus} \quad \@specOwner}
  \end{center}
  \vspace{0.4em}
  \hrule
  \vspace{0.8em}
}

\pagestyle{fancy}
\fancyhf{}
\fancyhead[L]{\small\color{specGray}\@specID}
\fancyhead[R]{\small\color{specGray}\@specTitle}
\fancyfoot[C]{\small\color{specGray}\thepage}
\renewcommand{\headrulewidth}{0pt}
\makeatother

\newtcolorbox{invariant}{
  colback=specBlue!5, colframe=specBlue!75,
  title=Invariant, fonttitle=\bfseries, breakable,
  boxrule=0.5pt, arc=2pt
}

\newtcolorbox{scenario}[1]{
  colback=white, colframe=specBlue!50,
  title={Scenario: #1}, fonttitle=\bfseries, breakable,
  boxrule=0.5pt, arc=2pt
}

\newtcolorbox{ac}[1]{
  colback=specAccent!5, colframe=specAccent!60,
  title={Acceptance Criterion #1}, fonttitle=\bfseries, breakable,
  boxrule=0.5pt, arc=2pt
}

\newtcolorbox{nongoal}{
  colback=specGray!8, colframe=specGray!60,
  title=Non-Goal, fonttitle=\bfseries, breakable,
  boxrule=0.5pt, arc=2pt
}

\lstdefinestyle{codestyle}{
  basicstyle=\ttfamily\small,
  breaklines=true,
  showstringspaces=false,
  frame=single,
  framesep=4pt,
  rulecolor=\color{specBlue!30},
  backgroundcolor=\color{specBlue!2},
  xleftmargin=0.5em,
  xrightmargin=0.5em,
  keywordstyle=\color{specBlue}\bfseries,
  commentstyle=\color{specGray}\itshape,
  stringstyle=\color{specGreen}
}
\lstset{
  inputencoding=utf8,
  extendedchars=true,
  literate=
    {—}{{--}}1
    {–}{{-}}1
    {…}{{...}}1
    {’}{{'}}1
    {‘}{{'}}1
    {“}{{``}}1
    {”}{{''}}1
    {é}{{\'e}}1
    {è}{{\`e}}1
    {ê}{{\^e}}1
    {à}{{\`a}}1
    {â}{{\^a}}1
    {ç}{{\c{c}}}1
    {ô}{{\^o}}1
    {→}{{$\to$}}1
    {←}{{$\leftarrow$}}1
    {×}{{$\times$}}1
    {≥}{{$\geq$}}1
    {≤}{{$\leq$}}1
    {≠}{{$\neq$}}1
    {∈}{{$\in$}}1
    {⌈}{{$\lceil$}}1
    {⌉}{{$\rceil$}}1
    {∞}{{$\infty$}}1
    {α}{{$\alpha$}}1
    {β}{{$\beta$}}1
}
\lstdefinelanguage{TypeScript}{
  keywords={const,let,var,function,return,if,else,for,while,class,interface,type,extends,implements,import,export,from,as,async,await,new,this,true,false,null,undefined,void,any,unknown,never,string,number,boolean,readonly,private,public,protected,enum},
  sensitive=true,
  comment=[l]{//},
  morecomment=[s]{/*}{*/},
  morestring=[b]",
  morestring=[b]',
  morestring=[b]`
}
\lstset{style=codestyle}

\newcommand{\Given}{\textbf{\color{specBlue}Given} }
\newcommand{\When}{\textbf{\color{specBlue}When} }
\newcommand{\Then}{\textbf{\color{specBlue}Then} }
\providecommand{\And}{}
\renewcommand{\And}{\textbf{\color{specBlue}And} }
 (adjust depth).

\documentclass[11pt,a4paper]{article}

\usepackage[utf8]{inputenc}
\usepackage[T1]{fontenc}
\usepackage{lmodern}
\usepackage[a4paper,margin=2.3cm]{geometry}
\usepackage{parskip}
\usepackage{microtype}
\usepackage[dvipsnames]{xcolor}
\usepackage{enumitem}
\usepackage{listings}
\usepackage{tcolorbox}
\usepackage{titlesec}
\usepackage{fancyhdr}
\usepackage{array}
\usepackage{longtable}
\usepackage{booktabs}
\usepackage{amsmath}
\usepackage{amssymb}
\usepackage{hyperref}
\usepackage{cleveref}

\tcbuselibrary{skins,breakable}

\definecolor{specBlue}{RGB}{30,64,132}
\definecolor{specGray}{RGB}{90,90,90}
\definecolor{specAccent}{RGB}{198,40,40}
\definecolor{specGreen}{RGB}{30,110,60}

\hypersetup{
  colorlinks=true,
  linkcolor=specBlue,
  urlcolor=specBlue,
  citecolor=specBlue,
  pdfborder={0 0 0}
}

\titleformat{\section}{\Large\bfseries\color{specBlue}}{\thesection}{0.75em}{}
\titleformat{\subsection}{\large\bfseries\color{specBlue!80}}{\thesubsection}{0.6em}{}

\makeatletter
\newcommand{\specID}[1]{\def\@specID{#1}}
\newcommand{\specTitle}[1]{\def\@specTitle{#1}}
\newcommand{\specStatus}[1]{\def\@specStatus{#1}}
\newcommand{\specVersion}[1]{\def\@specVersion{#1}}
\newcommand{\specOwner}[1]{\def\@specOwner{#1}}
\newcommand{\specTraces}[1]{\def\@specTraces{#1}}

\specID{SPEC-XXX}
\specTitle{Untitled Spec}
\specStatus{DRAFT}
\specVersion{0.1}
\specOwner{Jeremie Nunez}
\specTraces{}

\newcommand{\makespecheader}{%
  \begin{center}
    {\LARGE\bfseries\color{specBlue}\@specTitle}\\[0.4em]
    {\small\color{specGray}\texttt{\@specID} \quad v\@specVersion \quad \textsc{\@specStatus} \quad \@specOwner}
  \end{center}
  \vspace{0.4em}
  \hrule
  \vspace{0.8em}
}

\pagestyle{fancy}
\fancyhf{}
\fancyhead[L]{\small\color{specGray}\@specID}
\fancyhead[R]{\small\color{specGray}\@specTitle}
\fancyfoot[C]{\small\color{specGray}\thepage}
\renewcommand{\headrulewidth}{0pt}
\makeatother

\newtcolorbox{invariant}{
  colback=specBlue!5, colframe=specBlue!75,
  title=Invariant, fonttitle=\bfseries, breakable,
  boxrule=0.5pt, arc=2pt
}

\newtcolorbox{scenario}[1]{
  colback=white, colframe=specBlue!50,
  title={Scenario: #1}, fonttitle=\bfseries, breakable,
  boxrule=0.5pt, arc=2pt
}

\newtcolorbox{ac}[1]{
  colback=specAccent!5, colframe=specAccent!60,
  title={Acceptance Criterion #1}, fonttitle=\bfseries, breakable,
  boxrule=0.5pt, arc=2pt
}

\newtcolorbox{nongoal}{
  colback=specGray!8, colframe=specGray!60,
  title=Non-Goal, fonttitle=\bfseries, breakable,
  boxrule=0.5pt, arc=2pt
}

\lstdefinestyle{codestyle}{
  basicstyle=\ttfamily\small,
  breaklines=true,
  showstringspaces=false,
  frame=single,
  framesep=4pt,
  rulecolor=\color{specBlue!30},
  backgroundcolor=\color{specBlue!2},
  xleftmargin=0.5em,
  xrightmargin=0.5em,
  keywordstyle=\color{specBlue}\bfseries,
  commentstyle=\color{specGray}\itshape,
  stringstyle=\color{specGreen}
}
\lstset{
  inputencoding=utf8,
  extendedchars=true,
  literate=
    {—}{{--}}1
    {–}{{-}}1
    {…}{{...}}1
    {’}{{'}}1
    {‘}{{'}}1
    {“}{{``}}1
    {”}{{''}}1
    {é}{{\'e}}1
    {è}{{\`e}}1
    {ê}{{\^e}}1
    {à}{{\`a}}1
    {â}{{\^a}}1
    {ç}{{\c{c}}}1
    {ô}{{\^o}}1
    {→}{{$\to$}}1
    {←}{{$\leftarrow$}}1
    {×}{{$\times$}}1
    {≥}{{$\geq$}}1
    {≤}{{$\leq$}}1
    {≠}{{$\neq$}}1
    {∈}{{$\in$}}1
    {⌈}{{$\lceil$}}1
    {⌉}{{$\rceil$}}1
    {∞}{{$\infty$}}1
    {α}{{$\alpha$}}1
    {β}{{$\beta$}}1
}
\lstdefinelanguage{TypeScript}{
  keywords={const,let,var,function,return,if,else,for,while,class,interface,type,extends,implements,import,export,from,as,async,await,new,this,true,false,null,undefined,void,any,unknown,never,string,number,boolean,readonly,private,public,protected,enum},
  sensitive=true,
  comment=[l]{//},
  morecomment=[s]{/*}{*/},
  morestring=[b]",
  morestring=[b]',
  morestring=[b]`
}
\lstset{style=codestyle}

\newcommand{\Given}{\textbf{\color{specBlue}Given} }
\newcommand{\When}{\textbf{\color{specBlue}When} }
\newcommand{\Then}{\textbf{\color{specBlue}Then} }
\providecommand{\And}{}
\renewcommand{\And}{\textbf{\color{specBlue}And} }
 (adjust depth).

\documentclass[11pt,a4paper]{article}

\usepackage[utf8]{inputenc}
\usepackage[T1]{fontenc}
\usepackage{lmodern}
\usepackage[a4paper,margin=2.3cm]{geometry}
\usepackage{parskip}
\usepackage{microtype}
\usepackage[dvipsnames]{xcolor}
\usepackage{enumitem}
\usepackage{listings}
\usepackage{tcolorbox}
\usepackage{titlesec}
\usepackage{fancyhdr}
\usepackage{array}
\usepackage{longtable}
\usepackage{booktabs}
\usepackage{amsmath}
\usepackage{amssymb}
\usepackage{hyperref}
\usepackage{cleveref}

\tcbuselibrary{skins,breakable}

\definecolor{specBlue}{RGB}{30,64,132}
\definecolor{specGray}{RGB}{90,90,90}
\definecolor{specAccent}{RGB}{198,40,40}
\definecolor{specGreen}{RGB}{30,110,60}

\hypersetup{
  colorlinks=true,
  linkcolor=specBlue,
  urlcolor=specBlue,
  citecolor=specBlue,
  pdfborder={0 0 0}
}

\titleformat{\section}{\Large\bfseries\color{specBlue}}{\thesection}{0.75em}{}
\titleformat{\subsection}{\large\bfseries\color{specBlue!80}}{\thesubsection}{0.6em}{}

\makeatletter
\newcommand{\specID}[1]{\def\@specID{#1}}
\newcommand{\specTitle}[1]{\def\@specTitle{#1}}
\newcommand{\specStatus}[1]{\def\@specStatus{#1}}
\newcommand{\specVersion}[1]{\def\@specVersion{#1}}
\newcommand{\specOwner}[1]{\def\@specOwner{#1}}
\newcommand{\specTraces}[1]{\def\@specTraces{#1}}

\specID{SPEC-XXX}
\specTitle{Untitled Spec}
\specStatus{DRAFT}
\specVersion{0.1}
\specOwner{Jeremie Nunez}
\specTraces{}

\newcommand{\makespecheader}{%
  \begin{center}
    {\LARGE\bfseries\color{specBlue}\@specTitle}\\[0.4em]
    {\small\color{specGray}\texttt{\@specID} \quad v\@specVersion \quad \textsc{\@specStatus} \quad \@specOwner}
  \end{center}
  \vspace{0.4em}
  \hrule
  \vspace{0.8em}
}

\pagestyle{fancy}
\fancyhf{}
\fancyhead[L]{\small\color{specGray}\@specID}
\fancyhead[R]{\small\color{specGray}\@specTitle}
\fancyfoot[C]{\small\color{specGray}\thepage}
\renewcommand{\headrulewidth}{0pt}
\makeatother

\newtcolorbox{invariant}{
  colback=specBlue!5, colframe=specBlue!75,
  title=Invariant, fonttitle=\bfseries, breakable,
  boxrule=0.5pt, arc=2pt
}

\newtcolorbox{scenario}[1]{
  colback=white, colframe=specBlue!50,
  title={Scenario: #1}, fonttitle=\bfseries, breakable,
  boxrule=0.5pt, arc=2pt
}

\newtcolorbox{ac}[1]{
  colback=specAccent!5, colframe=specAccent!60,
  title={Acceptance Criterion #1}, fonttitle=\bfseries, breakable,
  boxrule=0.5pt, arc=2pt
}

\newtcolorbox{nongoal}{
  colback=specGray!8, colframe=specGray!60,
  title=Non-Goal, fonttitle=\bfseries, breakable,
  boxrule=0.5pt, arc=2pt
}

\lstdefinestyle{codestyle}{
  basicstyle=\ttfamily\small,
  breaklines=true,
  showstringspaces=false,
  frame=single,
  framesep=4pt,
  rulecolor=\color{specBlue!30},
  backgroundcolor=\color{specBlue!2},
  xleftmargin=0.5em,
  xrightmargin=0.5em,
  keywordstyle=\color{specBlue}\bfseries,
  commentstyle=\color{specGray}\itshape,
  stringstyle=\color{specGreen}
}
\lstset{
  inputencoding=utf8,
  extendedchars=true,
  literate=
    {—}{{--}}1
    {–}{{-}}1
    {…}{{...}}1
    {’}{{'}}1
    {‘}{{'}}1
    {“}{{``}}1
    {”}{{''}}1
    {é}{{\'e}}1
    {è}{{\`e}}1
    {ê}{{\^e}}1
    {à}{{\`a}}1
    {â}{{\^a}}1
    {ç}{{\c{c}}}1
    {ô}{{\^o}}1
    {→}{{$\to$}}1
    {←}{{$\leftarrow$}}1
    {×}{{$\times$}}1
    {≥}{{$\geq$}}1
    {≤}{{$\leq$}}1
    {≠}{{$\neq$}}1
    {∈}{{$\in$}}1
    {⌈}{{$\lceil$}}1
    {⌉}{{$\rceil$}}1
    {∞}{{$\infty$}}1
    {α}{{$\alpha$}}1
    {β}{{$\beta$}}1
}
\lstdefinelanguage{TypeScript}{
  keywords={const,let,var,function,return,if,else,for,while,class,interface,type,extends,implements,import,export,from,as,async,await,new,this,true,false,null,undefined,void,any,unknown,never,string,number,boolean,readonly,private,public,protected,enum},
  sensitive=true,
  comment=[l]{//},
  morecomment=[s]{/*}{*/},
  morestring=[b]",
  morestring=[b]',
  morestring=[b]`
}
\lstset{style=codestyle}

\newcommand{\Given}{\textbf{\color{specBlue}Given} }
\newcommand{\When}{\textbf{\color{specBlue}When} }
\newcommand{\Then}{\textbf{\color{specBlue}Then} }
\providecommand{\And}{}
\renewcommand{\And}{\textbf{\color{specBlue}And} }


\specID{SPEC-DB-002}
\specTitle{Typed Repositories --- Drizzle-inferred contracts for data access}
\specStatus{DRAFT}
\specVersion{0.1}
\specOwner{Jeremie Nunez}

\begin{document}
\makespecheader

\section{Context}
Repositories are the only code in the system allowed to issue SQL. They sit
directly above \texttt{@interview/db-schema} and expose typed, intention-named
functions to services. Thalamus cortices and the Sweep pipeline both consume
these repositories; keeping them thin, typed, and query-only is what lets the
service layer stay pure and testable.

\section{Goals}
\begin{itemize}[leftmargin=1.2em]
  \item Define the repository contract: \emph{queries only}, no business logic,
        no cross-entity aggregation beyond SQL joins.
  \item Parameter and return types come from Drizzle's
        \texttt{\$inferSelect / \$inferInsert}; nothing hand-rolled.
  \item Ban \texttt{any} and \texttt{unknown} in repository signatures. Payload
        narrowing happens above the repo, not inside it, and not at its edge.
  \item Make each repository method easy to mock at the service boundary.
\end{itemize}

\section{Non-Goals}
\begin{nongoal}
Repositories do not orchestrate multi-step workflows, call external services,
emit events, implement retry policies, or enforce business rules (tier gating,
permission checks, quota). They do not serialize to DTOs. That is the service
layer's job --- "repo = queries, service = logic". Transaction composition
spanning multiple repositories is owned by the service.
\end{nongoal}

\section{Invariants}
\begin{invariant}
No \texttt{any}, no \texttt{unknown}, no \texttt{Record<string, unknown>} in
repository public signatures. Parameters and return types are either Drizzle-
inferred row types, narrow DTO-less tuples, or primitives.
\end{invariant}

\begin{invariant}
A repository method performs at most one logical database operation: a single
\texttt{select} (with joins allowed), a single \texttt{insert}, a single
\texttt{update}, or a single \texttt{delete}. Composition of operations lives
in the service.
\end{invariant}

\begin{invariant}
Repositories take the \texttt{Database} handle (or a \texttt{tx}) by
constructor injection. They do not import a global client. This keeps them
trivially testable and transaction-safe.
\end{invariant}

\begin{invariant}
JSONB columns are not narrowed inside the repository. They are returned with
their schema-declared type; services are responsible for parsing via Zod or
equivalent before using them as domain objects.
\end{invariant}

\section{Interface}

\subsection{Construction}

\begin{lstlisting}[language=TypeScript]
import type { Database } from "@interview/db-schema";

export class ResearchFindingRepo {
  constructor(private readonly db: Database) {}
  // ... methods below
}
\end{lstlisting}

\subsection{Representative signatures}

\begin{lstlisting}[language=TypeScript]
import {
  researchCycle, researchFinding, researchEdge,
  type ResearchCycle, type NewResearchCycle,
  type ResearchFinding, type NewResearchFinding,
  type ResearchEdge, type NewResearchEdge,
} from "@interview/db-schema";

export class ResearchCycleRepo {
  constructor(private readonly db: Database) {}

  findById(id: bigint): Promise<ResearchCycle | null>;
  listByStatus(status: string, limit: number): Promise<ResearchCycle[]>;
  insert(row: NewResearchCycle): Promise<ResearchCycle>;
  updateStatus(id: bigint, status: string): Promise<void>;
}

export class ResearchFindingRepo {
  constructor(private readonly db: Database) {}

  findByEntity(
    entityName: string,
    entityType: string,
  ): Promise<ResearchFinding[]>;

  findByCycle(cycleId: bigint): Promise<ResearchFinding[]>;
  insertMany(rows: NewResearchFinding[]): Promise<void>;
}

export class ResearchEdgeRepo {
  constructor(private readonly db: Database) {}

  neighbors(
    name: string,
    type: string,
  ): Promise<ResearchEdge[]>;

  insert(row: NewResearchEdge): Promise<ResearchEdge>;
}
\end{lstlisting}

\subsection{What is explicitly forbidden}

\begin{lstlisting}[language=TypeScript]
// FORBIDDEN: leaks `any`
findSomething(input: any): Promise<any>;

// FORBIDDEN: leaks `unknown` payloads at the edge
getMeta(id: bigint): Promise<unknown>;

// FORBIDDEN: business logic in the repo
async publishArticleAndNotify(id: bigint) { /* no */ }

// FORBIDDEN: JS-side aggregation when SQL can do it
async listFindingsThenCountByType(cycleId: bigint) {
  const rows = await this.db.select()...;
  return rows.reduce(/* no -- push to SQL */);
}
\end{lstlisting}

\section{Scenarios}

\begin{scenario}{Service composes two repository calls in one transaction}
\Given \texttt{ResearchCycleRepo} and \texttt{ResearchFindingRepo} \\
\When the service wraps \texttt{updateStatus} + \texttt{insertMany} inside
\texttt{db.transaction(async (tx) => ...)} \\
\Then both repos accept the \texttt{tx} handle as their \texttt{Database} \\
\And a failure in either call rolls back the whole transaction.
\end{scenario}

\begin{scenario}{Lookup uses the declared index}
\Given \texttt{idx\_research\_finding\_entity} from SPEC-DB-001 exists \\
\When \texttt{ResearchFindingRepo.findByEntity("X","item")} is called \\
\Then the generated SQL filters on \texttt{(entity\_name, entity\_type)} \\
\And \texttt{EXPLAIN} reports an index scan, not a sequential scan.
\end{scenario}

\begin{scenario}{Unknown payload stays opaque at the boundary}
\Given \texttt{ResearchCycle.plan} is \texttt{jsonb} \\
\When a caller reads \texttt{cycle.plan} \\
\Then the type is the schema-declared type (not \texttt{any}) \\
\And the service --- not the repo --- validates and narrows it before use.
\end{scenario}

\begin{scenario}{Bulk insert preserves row shape}
\Given \(n\) \texttt{NewResearchFinding} values \\
\When \texttt{insertMany(rows)} is called \\
\Then a single SQL \texttt{INSERT} is issued \\
\And TypeScript rejects any row missing a \texttt{notNull} column.
\end{scenario}

\begin{scenario}{Mocked repo in a service unit test}
\Given a service under test whose only dependency is \texttt{ResearchFindingRepo} \\
\When the test substitutes a hand-rolled mock conforming to the repo's public
type \\
\Then no database connection is required \\
\And TypeScript enforces the mock implements every public method.
\end{scenario}

\section{Acceptance Criteria}

\begin{ac}{AC-1}
A lint rule (or grep-based test) finds zero \texttt{any}, \texttt{unknown}, or
\texttt{Record<string, unknown>} in exported repository signatures under
\texttt{packages/db-schema/src/repos/}.
\end{ac}

\begin{ac}{AC-2}
Every repository method's parameter and return types reference
\texttt{\$inferSelect} or \texttt{\$inferInsert} from
\texttt{@interview/db-schema}, or primitives (\texttt{bigint},
\texttt{string}, \texttt{number}), or arrays of those.
\end{ac}

\begin{ac}{AC-3}
No repository method contains control flow beyond a single await of a Drizzle
query builder expression. A static check (AST rule) enforces one logical
operation per method.
\end{ac}

\begin{ac}{AC-4}
Repositories accept a \texttt{Database | Transaction} handle at construction.
A test composes two repo calls inside a transaction and asserts rollback on
failure.
\end{ac}

\begin{ac}{AC-5}
\texttt{ResearchFindingRepo.findByEntity} uses
\texttt{idx\_research\_finding\_entity}; an integration test runs
\texttt{EXPLAIN} and asserts an index scan.
\end{ac}

\begin{ac}{AC-6}
A service unit test mocks every repo method from a single TypeScript
interface derived from the concrete class --- and compiles.
\end{ac}

\begin{ac}{AC-7}
Attempting to insert a row missing a \texttt{notNull} column fails at compile
time in a fixture test.
\end{ac}

\begin{ac}{AC-8}
No repository imports from \texttt{../services/}, \texttt{../controllers/}, or
any module outside \texttt{@interview/db-schema}. Enforced by an import-path
lint rule.
\end{ac}

\section{Traceability}

\begin{center}
\small
\begin{tabular}{@{}lll@{}}
\toprule
AC & Test file (planned) & Test name \\
\midrule
AC-1 & \texttt{packages/db-schema/tests/repos/no-any.spec.ts}         & \texttt{repo signatures ban any/unknown} \\
AC-2 & \texttt{packages/db-schema/tests/repos/types.spec.ts}          & \texttt{repo signatures use inferred types} \\
AC-3 & \texttt{packages/db-schema/tests/repos/single-op.spec.ts}      & \texttt{one logical query per method} \\
AC-4 & \texttt{packages/db-schema/tests/repos/tx.int.spec.ts}         & \texttt{rollback across two repos} \\
AC-5 & \texttt{packages/db-schema/tests/repos/finding-index.int.spec.ts} & \texttt{findByEntity uses entity index} \\
AC-6 & \texttt{packages/db-schema/tests/repos/mockability.spec.ts}    & \texttt{repo is mockable from its class type} \\
AC-7 & \texttt{packages/db-schema/tests/repos/insert-required.spec.ts}& \texttt{missing notNull column is compile error} \\
AC-8 & \texttt{packages/db-schema/tests/repos/isolation.spec.ts}      & \texttt{repos have no upward imports} \\
\bottomrule
\end{tabular}
\end{center}

\section{Open Questions}
\begin{itemize}[leftmargin=1.2em]
  \item Should repositories expose a cursor/stream API for large result sets
        (e.g. all findings for a long-running cycle), or is pagination enough?
  \item Do we introduce a \texttt{Repo} base class/interface for consistent
        \texttt{tx} propagation, or keep each class free-standing?
  \item Should JSONB columns be narrowed at the repo boundary via a branded
        type (still no \texttt{any}), or left to the service with Zod?
  \item How do we express "upsert" semantics without violating the
        one-operation-per-method invariant --- dedicated method, or service
        composition of \texttt{findById} + \texttt{insert}/\texttt{update}?
\end{itemize}

\end{document}
